% !TEX program = xelatex
% Cards format — auto-generated from comp-strats.tex by cards-convert.mjs
\documentclass[11pt]{article}

\usepackage{fontspec}
\usepackage{unicode-math}
% Latin Modern by OTF filename (not font name): xelatex on macOS (Core
% Text, no fontconfig) cannot resolve "Latin Modern Roman" by name, which
% silently produced ~4KB broken card PDFs. kpathsea finds the OTFs in
% texmf-dist on both macOS and the oven's Linux xelatex.
\setmainfont{lmroman10-regular.otf}[
  BoldFont=lmroman10-bold.otf,
  ItalicFont=lmroman10-italic.otf,
  BoldItalicFont=lmroman10-bolditalic.otf,
]
\setsansfont{lmsans10-regular.otf}[
  BoldFont=lmsans10-bold.otf,
  ItalicFont=lmsans10-oblique.otf,
  BoldItalicFont=lmsans10-boldoblique.otf,
]
\setmonofont{lmmono10-regular.otf}[Scale=0.88,
  BoldFont=lmmonolt10-bold.otf,
  ItalicFont=lmmono10-italic.otf,
]


\usepackage{graphicx}
\graphicspath{{figures/}}
\usepackage{booktabs}
\usepackage{tabularx}
\usepackage{adjustbox}
\usepackage{ragged2e}
\usepackage{microtype}
\usepackage{natbib}



\makeatletter
\def\input@path{{../}}
\makeatother
\usepackage{ac-paper-cards}


\hypersetup{
  pdftitle={Comp Strats},
}

\renewcommand{\acpdfbase}{comp-strats-26-arxiv}
\begin{document}

% ============================================================
% TITLE CARD
% ============================================================
\thispagestyle{empty}
\vspace*{0.3em}
\begin{center}
\includegraphics[width=0.62\textwidth]{figures/cover}\par\vspace{0.3em}
{\acbold\fontsize{18pt}{22pt}\selectfont\color{acdark} Comp Strats}\par
\vspace{0.1em}
{\fontsize{9pt}{11pt}\selectfont\color{acpink} Compositional Strategies in \ac{} --- Notes Toward a Theory of a Small Media-Arts Practice}\par
\vspace{0.4em}
{\normalsize\color{cyan!70!blue}\href{https://prompt.ac/@jeffrey}{\textbf{@jeffrey}}}\par
{\small\color{acgray} Aesthetic.Computer}\par
{\small\color{acgray} ORCID: \href{https://orcid.org/0009-0007-4460-4913}{0009-0007-4460-4913}}\par
\vspace{0.4em}
\rule{0.5\textwidth}{0.5pt}\par
\vspace{0.15em}
\colorbox{yellow!60}{\small\color{red!80!black}\textbf{\textit{working draft --- not for citation}}}\par
\vspace{0.1em}
{\footnotesize\color{acgray} March 2026 · \href{https://github.com/whistlegraph/aesthetic-computer/commit/b3a48ec6a}{b3a48ec6a}}\par
\end{center}

% ============================================================
% INDEX CARD
% ============================================================
\cardindex

% ============================================================
% BODY
% ============================================================
\section{The Condition}

A small media-arts practice operates under conditions that are not ordinary software conditions. A normal software project plans for growth, accumulates contributors, anticipates a customer base, and treats infrastructure as a resource to be scaled. A creative-software practice, by contrast, is a single author's instrument played over a span of years. Growth is suspect. Contributor accumulation is rare and often unwanted. Infrastructure is a hazard---and the hazard is not, in the first place, that infrastructure will fail. The hazard is that infrastructure will succeed, and in succeeding will outgrow the practice that motivated it. The practice ends up working for the platform rather than the other way around.

Five failure modes recur and are worth naming precisely. \emph{Capture}: a larger entity---a foundation, a platform company, an academic program---acquires the work or the maker and reframes both inside its own terms. The practice continues to exist on paper but stops being itself. \emph{Dilution}: the practice grows by accretion until its specifics dissolve into a brand or a vibe. \emph{Anthologization}: a curator, editor, or historian flattens a body of work into a representative sample inside someone else's narrative arc. \emph{Infrastructure rot}: the platform that hosted the work decays, and the work decays with it; the artifact is gone before its public can find it. \emph{Quiet abandonment}: the maker stops returning. Of the five, only the last is openly named, because only the last is sometimes willed; the others arrive uninvited, and are usually noticed too late to be defended against.

This paper asks: what strategies, internal to a small creative-software practice, are responsible for resisting these failure modes? Not by argument, not by manifesto, but in fact---visible in the choices the practice has actually made over six years. The answer is more interesting than a list of features. The strategies are not local optimizations. They are constitutive: they define what counts as a piece, an author, a composition, a refusal. They are also load-bearing in the engineering sense---removing any one of them produces drift, and drift is the surface form of every failure mode in the list above.

The paper proceeds analytically. Each section names a strategy or cluster of strategies, traces its sources where they are knowable, examines its costs as well as its payoffs, and surfaces the tensions it creates with other strategies. The strategies are not consistent with each other. The practice maintains them anyway, because consistency is a property of theories, and a working practice is not a theory. \citet{shklovsky1917art} called the operative technique of art \emph{defamiliarization}---making the familiar strange so that perception is forced where habit would otherwise let us sleepwalk. The strategies that follow are defamiliarization devices for the practice itself: each one names something the surrounding culture has rendered automatic, and refuses to let the practice render it automatic in turn. \citet{adorno1990punctuation} argued that even punctuation is constitutive, that small formal choices in notation shape what can be thought. The voice of the analysis below is the voice of the studio, not the conference~\citep{scudder2026voice}; the voice is itself a strategy under analysis.

\section{The Piece: A Theory of the Unit}

The unit of work in \ac{} is the \emph{piece}---a single \texttt{.mjs} or \texttt{.lisp} file that boots in a browser at a memorizable URL, runs a lifecycle of five verbs, and leaves~\citep{scudder2026pieces}. The decision to make the file the unit is consequential, and it is worth being precise about which decision it is. It is not the decision to make the URL the unit, although the URL co-evolved with the file~\citep{scudder2026url}. It is not the decision to make the user session the unit, the way SaaS products do. It is the file: a thing the author can hold in their head, edit in one buffer, and email if they had to.

The choice has antecedents. The Processing sketch is the closest cousin: a single \texttt{.pde} file that runs as a self-contained graphical program inside a development environment that does the boring work~\citep{montfort201210print}. The 1KB and 4KB demos of the European demoscene are an older cousin: a single binary that holds an entire animated piece, with every byte forced to do work because of the size constraint. Twine is another cousin: a single \texttt{.html} file that contains a complete hypertext story. In each case the unit is the smallest object that holds an entire piece of work. \ac{} inherits the form and converts it to its own platform; the constraint is the same constraint, the substrate is the browser, the tradition is shared.

The constraint costs something. A piece that does not fit in one file is forbidden by definition; if it does not fit, it is two pieces, or it is a piece plus a library, or it is a piece plus a session-server hook. The architecture of the platform separates piece-time concerns from system-time concerns precisely so that the piece itself can stay small, but the consequence is that some kinds of work are excluded: long simulations with extensive state, multi-piece games with cross-piece persistence, anything that wants to behave like a small operating system. Those exclusions are not bugs. They are what make the rest of the strategies possible. A unit that admits arbitrary complexity is not a unit; it is a project.

The constraint also pays. A single-file unit is a teachable unit, in the strict sense: an instructor can show one to a student, and the student can read it from top to bottom without losing the thread. A single-file unit is a coverable unit: another author can take the source and re-perform the piece without negotiating an entire repository. A single-file unit is an archivable unit: it can be mirrored, printed, embedded, or written to slow storage as a single object. The compression test that the platform applies---\emph{does this idea survive being written as a single file?}---is also the test of whether the idea is finished. Many ideas do not survive the test, and those ideas were not yet pieces.

The unit produces, not coincidentally, a particular cognitive economy. Programmers are accustomed to thinking in terms of \emph{what to change}: the diff is the unit of attention, and the existing code base is the substrate. Authors of pieces think in terms of \emph{what to draw}: the screen is the unit of attention, and the substrate is empty until they paint it. The shift from change-thinking to render-thinking has consequences elsewhere in the practice---most directly in the immediate-mode rendering model, where the screen is rebuilt from scratch every frame, and in the relationship between author and piece, where the piece is rehearsed rather than maintained. The single-file constraint is, in Adorno's sense, a punctuation mark for the practice: it is what stops a sentence from running on, and it determines what kind of sentences the practice can compose.

A productive open thread: pieces, even small ones, do not exist alone. They cluster---into series, into themed groups, into named collections that the author returns to over weeks or months. The unit between piece and platter is currently underspecified. The album, the program, the suite, the set, the cycle---all of these are candidate forms, all of them inherited from older arts. The practice has not yet committed to one. This is a productive gap rather than an oversight: the unit between piece and platter will be discovered by the practice, not designed in advance, and the analytic interest is in watching which inherited form turns out to fit.

\section{Time and Tempo}

Each piece runs a lifecycle of five verbs: \texttt{boot}, \texttt{paint}, \texttt{sim}, \texttt{act}, \texttt{leave}. The verbs are not arbitrary. They are an enactive model of an agent in a world: the agent boots into the world, perceives it (\texttt{paint} is misleadingly named, because in \ac{} \texttt{paint} is both the rendering pass and the perception cycle---the screen is what the agent shows the world it sees), updates internal state (\texttt{sim}), responds to inputs (\texttt{act}), and eventually departs (\texttt{leave}). The lineage is traceable to enactive cognitive science~\citep{varela1991embodied} and to the extended mind hypothesis~\citep{clark1998extended}: cognition is not a closed loop inside the head, it is the running of an agent in a world, with its instruments. The lifecycle codifies this for a piece, treating the piece as an agent and its author as an extended cognitive system that includes the piece.

The lifecycle has practical consequences. Anything that does not fit into one of the five verbs is suspect on inspection. Pieces routinely have boot-time setup, frame-rate paint loops, simulation steps that update at frame rate or otherwise, event-driven act handlers, and leave-time cleanup. They do not have init-deinit-render-update-handleEvent-onResize-onFocus-onBlur loops. The five verbs are a teaching device first and an architecture second; the architecture follows the teaching device, rather than the other way around. This inversion---teaching first, architecture second---is itself a strategy worth naming, because it determines which surface of the platform faces the new author.

The deeper choice in the lifecycle is the tempo at which the practice operates. A web service is \emph{deployed}; once deployed, it accumulates uptime, version numbers, customer contracts, and bug reports. A piece is \emph{practiced}: the relationship between author and piece resembles a musician's relationship with an instrument, not a developer's relationship with a service. Pieces accumulate hours, not commits. Hot reload is not a developer convenience but the equivalent of practicing at tempo, where the author plays the piece as it changes. The author and the piece are coupled in time, and the platform's job is to keep that coupling tight.

The cousins here are not other software practices; they are other instrument-and-rehearsal traditions. The live-coding tradition, in which code is improvised as a performance medium~\citep{magnusson2010designing}, is the closest. The musician's relationship with their instrument, rehearsing scales and figures over years until both are altered, is closer still. The classical-music score, where the score is the contract between the composer and any future performer, is more distant but informative: the score is a \emph{frozen} form of what \ac{} pieces try to keep \emph{live}. A piece in \ac{} sits between an instrument (which is rehearsed but not scored) and a score (which is frozen but not rehearsed); it is a thing the author keeps re-performing while it accumulates the form of having been written.

The hardest decision in the lifecycle is the last verb: \texttt{leave}. There is no merge-to-main, no v1.0, no production cutover. A piece is finished when the author stops returning. This is a deliberate refusal of the version-number contract that engineering inherits from manufacturing, and it opens a hard question: how does a piece accumulate longevity if not through versioning? The practice has partial answers. The URL persists. The source is downloadable. The piece can be covered by another author, in which case it returns to active practice through a different hand. But each answer is partial, and longevity in a non-versioned medium remains a real open thread.

A second open thread, from the same source: many pieces do not actually leave. They keep being returned to. The practice has produced no good vocabulary for this state, in which a piece is neither finished nor abandoned but living in a long tail of intermittent revision. Calling such a piece a \emph{draft} is wrong, because drafts imply a final form they are aiming at; calling it \emph{active} is also wrong, because active implies recent work. The right vocabulary is probably borrowed from oral tradition: pieces in a long tail are \emph{being told}---kept alive by re-performance, not by maintenance. If that is the right word, it has consequences for how the platform should treat such pieces---perhaps as a third state alongside published and unpublished, perhaps not at all. The thread is still open.

\section{Names and the Prompt}

\ac{} routes everything through the URL~\citep{scudder2026url}. A piece is \texttt{aesthetic.computer/notepat}. A user piece is \texttt{aesthetic.computer/@user/piece}. Parameters arrive after colons, in the order the piece declares them. The namespace is the interface: there is no menu, there is no dashboard, there is no \emph{browse} action. To use the platform, you have to know the name.

The strategy is to make the URL \emph{mnemonic} rather than \emph{addressing}. A mnemonic URL is a compressed path back to a practice. An addressing URL is a pointer to a resource. A user who has internalized \texttt{notepat}, \texttt{wrl}, \texttt{prompt}, \texttt{clock} has internalized a map of possible experiences, the way a regular at a venue has internalized which staff members work which shifts. The platform rewards memorization in the same way a city rewards memorization. There is no shortcut for someone who has not yet learned the streets.

The mnemonic strategy has obvious costs. Discovery is harder. The platform does not greet new users with a curated home page. The home page, such as it is, is the prompt itself, and the prompt assumes the user has somewhere to type. The cost is paid in onboarding, and it is mitigated only by the fact that \ac{} pieces tend to share through trusted channels---a friend, a class, a card, a paper---rather than through search, so the relevant first contact is rarely cold. The cost is also redirected: the practice spends in onboarding what it would otherwise spend in feed design, recommendation systems, and engagement loops. The redirection is intentional. Time spent designing a recommendation system is time not spent designing pieces.

The deeper strategy is the elevation of the prompt itself. The command line at the top of every screen is the platform's instrument, designed to reward fluency rather than configuration~\citep{scudder2026notepat}. A user who has internalized the prompt can move between pieces without ever touching a menu. The prompt is typographically ascendant in the interface---a design decision as much as an interaction one---because the practice's claim is that the typed name is the most powerful interaction primitive available in a browser, and the interface should make this visible.

The lineage of the prompt-as-instrument is long and largely unwritten in the contemporary literature. The Smalltalk workspace, the Lisp REPL, the Plan 9 \texttt{rc} shell, the HyperCard message box: each is a place where the user can type a small symbol and have a system act. The \ac{} prompt is closer to these than to the modern application's \emph{command palette}, which is generally a search interface over a finite set of features. The \ac{} prompt is open: a user can type the name of any piece, and the system will load it. Failure to find a name produces a piece, not an error. The distinction between a search palette over closed features and a typed instrument over an open namespace is the distinction between a touch-tone phone and a piano.

The prompt produces an aesthetic. Names of pieces in \ac{} tend to be short, evocative, two-syllable, lowercase: \texttt{notepat}, \texttt{prompt}, \texttt{clock}, \texttt{wrl}. The aesthetic is closer to Erik Satie's piano pieces or Aphex Twin's track titles than to the descriptive titles of typical software. The name is the first thing the audience repeats; if it is not repeatable, it is not yet named. The practice has produced an internal taste-test that operates almost without conscious deliberation: a candidate name that does not repeat well is dropped before it is published.

A productive tension: the namespace is finite. As the platform grows, good names get taken. The practice has not yet developed an explicit strategy for namespace stewardship beyond first-come-first-served, and the open thread is whether the platform should evolve toward a Twitter-style handle system (where the personal namespace is unbounded but the global namespace is short and contested) or toward something closer to a magazine's table of contents (where the global namespace is curated and finite, and entries are added over editorial cycles). Both have costs. A handle-style namespace privileges early adopters; a curated-style namespace privileges the curator. The practice will probably arrive at a hybrid, but the form is still unspecified.

\section{Authorship: Cover, Not Fork}

Authorship in \ac{} is musical. Pieces are \emph{covered}, not \emph{forked}. The platform's \texttt{source} command hands the requesting user the source of a named piece and authorizes them to publish it as their own work, signed with their handle. This is deliberately not a Git fork. A fork carries a parent--child relationship, a divergence narrative, and a long-running expectation of merge or reconciliation. A cover does not. A cover is a respectful re-interpretation; the original persists, the cover persists, and the relationship between them is the listener's to decide.

Stiegler's concept of the \emph{idiotext} is useful here~\citep{stiegler1998technics,staunaes2021stiegler}. The idiotext is the singular memory woven through a person's technical prostheses---their notebooks, their files, their instruments, their habits. The point of an authored piece is not, on this account, that it expresses the author. The point is that it remembers for them. A piece holds a fragment of the author's thinking that the author no longer has to hold. The piece is, in this sense, a prosthesis for memory. A cover, then, is a particular kind of remembering: one author taking another author's idiotext and replaying it through their own. The cover is not a tribute; it is an act of cognitive partnership, in which two memories share a substrate.

This produces a different metaphysics of authorship than software typically affords. In software, the question \emph{who wrote this?} admits one answer (the original committer), one footnote (the contributors list), and one legal answer (the license holder). In \ac{}, the question admits multiple parallel answers, each with its own URL, its own handle, its own moment of publication. The original is not threatened by covers. The covers are not derivative in the legal sense; they are derivative in the musical sense, which is the kind of derivation that produces new work. The voice in code, in commits, in documentation, and in papers is the same voice~\citep{scudder2026voice}, and it remains first-person throughout: \emph{I built this} beats \emph{the author implemented} because the latter pretends to a neutrality that no actual author has ever inhabited.

The cover model has an unsolved tension. \emph{Cover} is an apt metaphor when one author re-performs another author's piece. It is not a metaphor for collaboration on a single piece, in which two or more authors contribute simultaneously. The practice has so far avoided multi-author single pieces in the strict sense, but multi-player pieces---\texttt{squash}, \texttt{1v1}, \texttt{udp}---are a partial form of co-authorship in which the runtime is shared even when the source is not. The right vocabulary for runtime co-authorship is probably \emph{session} or \emph{duet} or \emph{improvisation}, each borrowed from music. None of these have yet been formalized in the platform, and the open question is whether they should be---whether multi-author pieces want to be a first-class kind of work, or whether the cover model is sufficient.

A second tension: publishing is performance. Issuing \texttt{publish} is not deployment; it is stepping on stage. The piece enters the namespace and becomes playable by anyone with the URL. There is no staging environment, no canary deploy, no feature flag. There is only the live URL. The model assumes---perhaps demands---that the author is rehearsed. A piece that is published unrehearsed is a public failure, in the same way that a musician who steps on stage unprepared is a public failure. The practice treats this as a feature: the consequence of unrehearsedness is artistic, not infrastructural, and is therefore survivable. Software's elaborate apparatus of staging, feature-flagging, and rollback exists to make unrehearsed work invisible to the audience. \ac{} does not provide that apparatus, on the principle that an unrehearsed piece should be visible as such, because visibility is what motivates rehearsal.

\section{Composition Proper}

The strategies that earn the paper its name concern how works are assembled from the platform's primitives. They cluster around five claims, each of which deserves more space than the heading-and-gloss form previously given to it.

The first claim is that \emph{composition begins from instruments, not from prompts}. A song made by typing a sentence into a closed-source generative endpoint is not a composition in the sense the practice cares about; it is a request. The endpoint composes; the user prompts. The result may be musical, may be moving, may be commercially successful, but it is not a piece in the \ac{} sense, because the piece is what the author wrote, and the user wrote a sentence. Pop research in \ac{} accordingly starts from \texttt{notepat}, \texttt{sinebells}, \texttt{chord}, \texttt{beat}---instruments authored inside the system---and builds upward. The pipeline composes from these primitives in the same way an early-1980s synth-pop record composed from a TR-808 and a DX7: each instrument is itself a piece, with its own voice and limits, and the composition is the negotiation among them. End-to-end generative composition is excluded by design, not because the results would be unmusical but because the results would not be inspectable, and inspectability is a load-bearing requirement everywhere else in the practice.

The second claim follows from the first: AI is a \emph{performance layer}, not a composer. A trained model can sing a part the same way a sampler can play a part---as one voice in an arrangement the human authored. The arrangement is the composition. The arrangement must survive being inspected. The analogy here is to the synthesizer in 1960s pop and the sampler in early hip-hop: both were initially treated as cheats by the surrounding culture, both turned out to be instruments, and both became compositional only when their authors learned to arrange them rather than rely on their factory voices. Vocal generation in the current pipeline routes through a third-party text-to-speech provider for performance reasons and is treated as an instrument with limits to be negotiated; an \ac{}-native vocal lane has been attempted and dropped, with the question of whether to resume it---perhaps via a vocal-tract physical model rather than a learned model---still open.

The third claim concerns legibility: \emph{the seam shows}. \ac{} does not paper over the boundary between a sinebells bed and a notepat lead, between a rendered subtitle and a live overlay, between a human-drawn line and a machine-assisted line. The seam is information. It tells the listener how the work was made. The listener is, on the practice's terms, a peer rather than a consumer; therefore the listener has a right to know how the work was made. A smoothed seam is forgery, in the literal sense: the signature of one process is altered to look like another. The visible-seam policy is stronger than the comparable policy in fine-art photography (where retouching is concealed but expected) or in commercial popular music (where pitch correction is concealed but expected). It is closer to disclosure policies in scholarly publishing (where methods sections are required) or to early-modern oil painting (where pentimenti---visible traces of overpainting---are read as evidence of process and valued accordingly). The policy is also a defense: a smoothed seam is the surface form of capture, in which a work is reframed inside someone else's narrative, while an unsmoothed seam resists capture by making the work's authorship inarguable.

The fourth claim concerns scoring: if you cannot write the score, you do not know the piece. The score is what makes a piece teachable. In \ac{}, the score takes various forms: the source file, the prompt log, the \texttt{waltz-events.json} file that drives the recap pipeline, the \texttt{papers/SCORE.md} index that drives the platter, the per-word \texttt{.edit.json} recipe that drives vocal post-production. The score is not a visualization of the piece; it is the piece's notation, the form in which the piece can be re-performed by another author or by a future version of the same author. The score discipline has a useful side effect: it forces the practice toward legible processes. A pipeline that depends on irreproducible inputs---a magic API key, a one-off model run, an undisclosed prompt---cannot be scored, and therefore cannot be a piece. This excludes a great deal of contemporary AI-driven art on principle. The practice does not exclude AI; it excludes processes that cannot be written down. The difference matters.

The fifth claim is the immediate-mode discipline: \texttt{paint} is called every frame, and the screen is rebuilt from scratch~\citep{montfort201210print}. There is no retained scene graph. The cognitive payoff is that the artist thinks in terms of \emph{what to draw now} rather than \emph{what to change}. The shift is the same shift that distinguishes drawing from sculpting, or watercolor from oil. Drawing, like immediate-mode rendering, is a discipline of the present frame; the previous frame is gone. Oil paint and retained scene graphs both allow longer-running negotiations with the surface, and their cost is a kind of temporal blur in which the artist is always working against an inherited state. Immediate-mode is the practice's commitment to drawing-time over editing-time, and the commitment is nearly absolute---retained state is a per-piece decision rather than a platform default.

A sixth claim runs across the others: \emph{translate, always}. Every paper exists in four languages~\citep{scudder2026score}, and the work is not finished until it has left English. The reasoning is partly civic, in that the practice does not believe English is the only language in which thinking happens; and partly defensive, in that a paper that can be read in only one language can be captured by the institutions of that language. Translation is a form of distribution that no single institution controls. It is also a re-test of the score: a paper that survives translation is a paper that has been read, by someone with no particular stake in defending its English version, and has held up.

\section{Refusal as Constitutive}

A practice is also defined by what it will not do. The refusals in \ac{} are non-negotiable, and they are constitutive in the strict sense: removing any one of them produces a different practice, not a less defended one. Negative space, in drawing, is the unmarked region that defines the marked one; in a practice, refusals do the same work. They give the practice a recognizable silhouette, and a practice without sharp refusals has no silhouette at all.

Closed source has no place in schools. The pedagogical argument is given at length elsewhere~\citep{scudder2026openschools}; the compositional argument is shorter. A black box cannot be taught. Anything taught must be openable by the student, because the student's authority over the piece is what makes the lesson a lesson rather than a demonstration. A closed-source piece in a classroom is, structurally, a vendor presentation; the teacher mediates between the vendor and the student, and the student's role is to receive. An open-source piece in the same classroom is, structurally, a workshop; the teacher and the student both have authority over the artifact, and the lesson proceeds by both modifying it. The two situations are not differently shaped versions of the same thing; they are different things.

No analytics, no telemetry, no dark patterns. The platform does not measure its users. It is not optimizing them. It is hosting them. The first behavior follows from a conviction about the audience: the audience is a peer, not a metric. A platform that measures its users is treating them as a population to be managed, and the practice rejects that framing on aesthetic as well as ethical grounds. The aesthetic ground is real and worth stating: measured users behave less interestingly than unmeasured ones, because they are aware they are being watched. The observer effect applies to interfaces. A practice that wants to see the audience as it is must decline to watch it.

Drafts only; the send is human. Outreach, applications, pitches, replies---automated assistance produces drafts, never sends. The human remains the only signature that has ever been on the work. The reasoning is partly about reliability (an automated send can hallucinate, can misaddress, can fail in ways the human review catches) and partly about politics (the human signature is the locus of accountability, and accountability cannot be automated). The send rule generalizes: anything that leaves the studio under the studio's name leaves under a recent and explicit human decision.

What the private archive keeps, the public repo does not surface. The studio has a vault---an encrypted, private archive of materials that have not been cleared for publication. The vault is not the platter. Material flows from vault to platter only after a human review, and the review is conservative. This is the archive's analogue of the send rule. It also does important defensive work: the vault contains material that, if surfaced prematurely, would change the practice's relationships with its collaborators, family, and audience. Conservative review is the price the practice pays for being able to keep certain things in the vault at all.

Never anonymize the labor. Every commit, every cover, every paper is signed. The anti-pattern---unattributed work flowing through a corporate channel---is what the field is trying to escape, not reproduce. A practice that tolerates unattributed labor cannot defend that labor against capture, because the capturing party will simply add their own attribution to a piece that has none. Attribution is therefore not a vanity; it is a load-bearing layer in the defense.

A productive tension across the refusals: each one closes off a class of revenue or scale that other small studios use to survive. Telemetry-driven product development is excluded; closed-source teaching licenses are excluded; automated-outreach growth is excluded; vault-to-public surfacing for engagement is excluded; anonymous labor (the form of labor most easily monetized at scale) is excluded. The refusals are expensive. The open question of \S11 is whether their cost is sustainable.

\section{Care as Method}

Care is a method in this practice, not a mood. Warmth is the default. The voice for collaborators, students, and family is warm before it is correct. Coldness is the failure mode of efficiency culture; in a small studio, coldness is also the failure mode of burnout. A warm voice maintained over years is also defensive infrastructure: it keeps the trusted channels open, and trusted channels are the ones through which opportunities arrive.

Opportunities that arrive through trusted channels are written down in a single inbox, the studio's \texttt{honeydo.txt}. They are not a backlog. They are a queue of relationships. The distinction matters: a backlog is a list of work that someone (probably the author) is supposed to do, while a queue of relationships is a list of human ties that are currently warm and require action---a reply, a draft, a meeting, a follow-up. Confusing the two collapses the practice into productivity software, in which relationships are reduced to tasks and lost on completion.

The dossier swimlane in the platter~\citep{scudder2026score} is care extended into the field. The dossiers record what is publicly recoverable about the organizations that fund and host media arts: financials, governance, programs, named grants, conflicts of interest. The work is fact-surfacing rather than argumentative; the dossier ends where the facts end, and does not draw conclusions the facts will not support. This posture is itself an act of care for the field: the field is small enough that the same names recur across contexts, and a clear factual map of the field is a kind of protective infrastructure for everyone working inside it. The dossier rule for code comments and commit messages is the same: surface the facts, stop where the facts stop.

A useful distinction, between care and customer service. Customer service responds to needs the customer has articulated; care anticipates needs the recipient has not articulated. The practice tries to be care, not service. A piece that anticipates a teacher's classroom need is care; a piece that responds to a teacher's filed feature request is service. Both are valuable. The practice is biased toward the first, partly because the first scales worse and is therefore rarer in the surrounding industry, and partly because the second tends to draw the practice into a feedback loop in which it begins to optimize for filed requests rather than for the unrequested needs that the audience cannot articulate but will recognize when met.

\section{The Platter}

The platter is the long-running collection from which papers, pieces, and dossiers are carved~\citep{scudder2026score}. It is the largest unit in the practice---larger than any single paper or piece, slower than any single deadline. Papers are added before they are finished. Threads accumulate without obligation to become papers. Dossiers expand as the field's facts change. The platter's job is to hold the raw material that the mill (the paper-generation pipeline) will eventually press.

The platter has rules, three of which are worth analyzing in turn. The first: if a thread can survive being written down, it was real. Writing is the test, not the documentation. Anything that disintegrates the moment the author tries to write it was scaffolding for the conversation, not signal for the field. The rule is brutal in practice: many ideas that felt important in conversation evaporate on the page. The platter accepts this as a feature. The page is the test. An idea that cannot survive the page may still be valuable in conversation---in fact, the spoken voice is often where the practice already does its sharpest distillation---but it is not yet a platter idea, and the platter resists letting un-tested speech occupy its space. This is the strict version of why the present paper has had to be rewritten away from its previous aphoristic form: aphorisms are the form of things that work in speech, and the platter is the place for the things speech alone cannot finish.

The second rule: cover what you cite. Every visual claim in a dossier cover has a source line; every reference in a paper points to a thing the author has read, not a thing they have heard of. Defensible footing is the only footing. The rule is restrictive: the practice cannot cite the things it has heard of but not read, and the things it has heard of but not read are often the things it would most like to cite. The discipline produces honest papers at the cost of slower ones. The trade is taken willingly, because the alternative---papers full of name-dropped references---would be a particular kind of forgery that the practice has elsewhere ruled out.

The third rule: mill, don't market. The papers in the platter exist because the practice generates more thinking than the practice can absorb. The mill catches the overflow. It is not a marketing pipeline; it is an honesty pipeline. The distinction matters: a marketing pipeline produces materials calibrated to a target audience, while an honesty pipeline produces materials calibrated to the practice itself. Audiences may or may not find the materials; that is their decision, not the mill's brief.

The platter has an unresolved structural question. A vault is private; a platter is public; what is the third state? Many materials in the practice---sketches, drafts, partial dossiers, half-cited threads---are not yet ready for the public platter, but are ready for trusted peer eyes. The practice has not yet built a third state, and the open thread is whether one is needed. A peer-only state could be implemented as a private mirror, as a password-protected URL, as a Zotero group, or as something natively \ac{}. The question is whether such a state would erode the simple two-state architecture (vault vs.\ platter) on which the send rule depends, or whether it would in fact strengthen the architecture by giving half-finished work a legitimate home that does not require either private encryption or public exposure.

\section{The Shield: A Theory of Layered Defense}

The strategies above do not exist for their own sake. They are organized to defend against the five failure modes named in \S1, and the defense is layered: aphorism, piece, paper, dossier, platter. (The label \emph{aphorism} is retained as a name for the smallest unit of teachable doctrine, even though this paper is no longer written in aphorisms; the argument for compression as a defense survives even when the surface form of the paper does not perform it.) Each layer slows a different mode of attrition.

Capture is resisted by the mobility of the strategies themselves. Strategies travel; institutions cannot acquire a sentence. A practice whose substance is encoded in transferable strategies---and not in proprietary platform features or in personal charisma---is harder to capture than one whose substance is concentrated in either, because the capturing party would have to capture every author who has internalized the strategies, and the strategies, once internalized, do not require the original author to keep functioning.

Dilution is resisted by the refusals (\S7). A practice grows by accretion until something stops it; the refusals stop it. They define what the practice is by defining what it is not. A practice that refuses closed source in schools, that refuses analytics, that refuses automated sends, that refuses to publish from the vault, that refuses to anonymize labor---that practice has a recognizable shape, and the shape resists being smoothed into an industry vibe.

Anthologization is resisted by the dossier posture and by the visible seam. A work that surfaces its own context and process makes anthologization harder, because the curator who tries to flatten the work into a sample finds the work already speaks for itself. The dossier is a counter-anthology: a record of context that travels with the work and arrives before the curator does.

Infrastructure rot is resisted by the architecture: pieces are single files at named URLs, the platter is the canonical archive, papers exist in four languages, and the platform's source is openable. Each redundancy is a hedge against a particular mode of decay. No single hedge is sufficient---a single-file piece is still vulnerable to platform rot, a translated paper is still vulnerable to library decay---but the layers compound. The bet is that no realistic mode of attrition attacks all layers at once.

Quiet abandonment is the hardest mode to defend against, because it is sometimes the right outcome. A piece that is done is sometimes a piece whose author has chosen to stop returning. The practice converts this from a failure mode into a positive value: \emph{done when abandoned}. The piece persists at its URL. The source is downloadable. Another author can cover the piece, in which case it returns to active practice through a different hand. The defense against quiet abandonment is not that the practice continues forever, but that the practice's substrate---the URL, the source, the cover convention, the platter---outlives any single author's attention.

\citet{magnusson2010designing} describes the epistemic dimension of musical instruments: instruments are tools that shape what their players can think. \ac{} is its own instrument, and the strategies in this paper are the technique. They are not optional, in the same way a violinist's bowing technique is not optional. The instrument is playable without them. It is not playable \emph{well} without them. The shield is the technique made explicit; the practice is the instrument played.

A note on the metaphor. \emph{Shield} is martial, and the metaphor warrants questioning. A martial metaphor implies an adversary; the failure modes named here are not adversaries in the conscious sense. (Capture is an outcome, not an actor; dilution is a process, not an attacker.) A more accurate metaphor might be ecological: the strategies are an immune system, or a coral reef, or a soil structure that keeps a small ecosystem alive in the presence of larger ecosystems that do not specifically wish it harm but do not specifically wish it well either. The ecological metaphor would emphasize redundancy, slow accumulation, and mutualism rather than defense per se. The practice has not yet committed to one metaphor over the other. The shield language is current; the soil language is plausible. Future versions of this paper may move between them.

\section{Open Threads}

A real shield has gaps. So does this one. Four gaps are large enough to name explicitly.

\emph{Economics.} The strategies above describe a posture. They do not describe how the studio pays rent. Sustainability is treated separately~\citep{scudder2026sustainability}; the synthesis between the compositional strategies in this paper and the economic conditions in that one is unfinished. There is real friction. The refusals (\S7) close off classes of revenue---telemetry-driven, attention-extracting, closed-source-licensing---that other small studios use to survive. The practice has not yet articulated a revenue model that is consistent with the refusals, and the open thread is whether such a model exists or whether the practice will have to be partly subsidized---by grants, by patronage, by adjacent labor---for as long as the refusals hold. The honest answer at present is the second.

\emph{Successor culture.} The strategies are currently held by one author. A culture is what survives the author. The translation, dossier, cover, and platter strategies are first attempts at building artifacts that survive any one author's attention, but the rest of the question is unsolved. The deepest version of the question is whether a single-author practice can produce a multi-author culture without going through an institutional phase. Most cultures that have done this have done it through schools, journals, or guilds. \ac{} has gestures toward each (the open-schools paper, the platter, the cover convention) but has not yet committed to one form. The thread is alive; the synthesis is open.

\emph{Adversarial robustness.} Strategies can be repeated by an actor who does not believe them. The defense against rhetorical theft is not in the words themselves but in the surrounding artifacts---the pieces, the papers, the platter. This is why the shield is layered. But the layered defense has a known weak point: an adversary who does not need to win, only to confuse, can degrade the practice's recognizability by repeating its sentences without its substance. The practice has not yet articulated a strategy against such an adversary. The open thread is whether such a strategy can be articulated at all, or whether the only available defense is the ongoing demonstration that the practice's sentences are backed by its work.

\emph{Scale.} The practice remains small by choice. If it grows---through covers, through schools, through pieces that escape the author's attention---the strategies above will need to be tested at scale. Several of them (the refusals, the cover convention, the dossier posture, the visible-seam policy) are likely to scale. Several (warmth-as-default, the trusted-channels theory of opportunity, the conservative vault review) will need to be reformulated. The open thread is which strategies are essentially small and which are essentially scale-invariant, and the analytic interest is in finding out which is which by trying. Some failures of scale will not be visible from inside the small practice; they will only be visible after the strategy has already broken under load.

\section{Conclusion}

Compositional strategy in \ac{} is not a feature list. It is a posture, encoded in strategies short enough to teach and analytical enough to extend. Each strategy has been examined here for what it presupposes, what it costs, what it buys, and where it tensions other strategies. The strategies are not consistent with each other, and they were not designed to be. They are load-bearing in a practice that is older than any of them individually, and the practice maintains them because removing any one of them produces drift.

The deeper claim is that the media arts deserve a doctrine---meaning a body of strategies that fits in the head of a single author and is teachable to the next. The doctrine has to be small enough to carry and large enough to be analyzed. This paper is the analytic side of the doctrine; the strategies it analyzes are, taken individually, smaller than the analysis. That is the right relation between the two: the strategies fit on a card, and the analysis explains why they belong on the card. Earlier drafts of this paper attempted to compress the analysis back into the card and lost most of what mattered. The aphoristic form turned out to be a poor compression of things that the spoken voice already says better in the studio; the analytic form is the form for the things speech alone cannot finish.

The shield is unfinished. The gaps are real. The practice is younger than the questions it is trying to answer. A finished shield is a finished practice, and a finished practice is one that has stopped finding new ideas. The platter is always growing, the doctrine is always being re-tested, and the analysis, like this paper, is partial on purpose. What the paper proposes is not a completed defense but the form of a defense that is willing to keep being written.

\bibliographystyle{plainnat}
\bibliography{references}

\end{document}
